
\subsection{Consideraciones y precauciones hardware}

Antes de energizar el banco de pruebas, se debe verificar que todas las conexiones de potencia y señal coincidan con los diagramas del montaje. En particular, es importante revisar la polaridad de las fuentes, la continuidad de las tierras de referencia y la correcta conexión de los terminales de armadura y campo de las máquinas de corriente continua.

La carga electrónica y las etapas de acondicionamiento deben operarse dentro de los rangos de tensión y corriente definidos durante el diseño. No se recomienda aumentar la corriente de carga sin supervisar la temperatura del transistor de potencia y del disipador, ya que una disipación excesiva puede dañar el componente o modificar el comportamiento del sistema durante los ensayos.

También se debe evitar manipular el cableado con el sistema energizado. Cualquier cambio en las conexiones del motor, generador, sensores o carga electrónica debe realizarse con las fuentes desconectadas y verificando posteriormente que no existan cortocircuitos entre alimentación, tierra y señales analógicas.

\subsection{Consideraciones para simulación en tiempo real con Simulink}
\label{subsec:simulink_consideraciones}

La ejecución del modelo en Simulink depende del rendimiento del computador y de la estabilidad de la comunicación Ethernet con el OPTO22. Por este motivo, se recomienda utilizar un tiempo de muestreo $T_s$ coherente con la capacidad real del sistema, evitando valores demasiado pequeños que puedan generar retrasos, pérdida de sincronismo o saturación del procesamiento.

Durante las pruebas, se debe limitar el uso de bloques de visualización como \emph{Scopes}, registros excesivos de datos u otros elementos que aumenten innecesariamente la carga computacional. Asimismo, se recomienda cerrar aplicaciones en segundo plano y mantener el computador dedicado exclusivamente a la ejecución del modelo mientras se realizan ensayos en tiempo real.

\subsection{Consideraciones y precauciones módulos OPTO22}
\label{subsec:opto22_consideraciones}

\subsubsection{Módulos de E/S utilizados}
El banco de pruebas se implementa con módulos SNAP de OPTO22:
\begin{itemize}
    \item \textbf{Dual Analog Output \texttt{SNAP-AOV-27}}: 2 canales de \emph{salida analógica de voltaje} bipolar, con rango nominal de \(-10\) a \(+10\)~VDC. Cada canal puede suministrar solamente \(\pm 5\)~mA de corriente de carga, por lo que \emph{no debe utilizarse como driver de potencia}. Ambos canales comparten referencia y \textbf{no están aislados entre sí}, aunque el módulo sí posee aislamiento respecto a otros módulos o dispositivos del rack. \cite{opto22_snap_aov27}
    \item \textbf{Quad Analog Input \texttt{SNAP-AIV-4}}: 4 canales de \emph{entrada analógica de voltaje}, configurables por canal en \(\pm 10\)~VDC o \(\pm 5\)~VDC. Los canales \textbf{no están aislados entre sí}. \cite{opto22_snap_aiv4}
\end{itemize}

\subsubsection{Límites eléctricos y protección del hardware OPTO-22}
Para evitar daños en los módulos y/o en la electrónica asociada, se deben respetar las siguientes restricciones:
\begin{itemize}
    \item \textbf{No exceder el rango de salida del \texttt{SNAP-AOV-27}}: asegurar que la señal de comando enviada desde Simulink se mantenga dentro de \([-10, +10]\)~V, idealmente con saturación explícita en el modelo.
    \item \textbf{No exigir corriente de carga mayor a \(\pm 5\)~mA por canal} en el \texttt{SNAP-AOV-27}. Si se requiere accionar una etapa de potencia, debe usarse un amplificador/buffer externo de alta impedancia de entrada (p.\,ej. etapa de acondicionamiento) para que el módulo no entregue corriente directa a la carga. 
    \item \textbf{Entradas del \texttt{SNAP-AIV-4}}: garantizar que las señales medidas no superen el rango configurado (\(\pm 10\)~V por canal).
    \item \textbf{Referencias/tierras}: dado que los canales no están aislados entre sí (entrada y salida), es crítico compartir referencias de señal de forma consistente y evitar lazos de tierra.
\end{itemize}

\subsubsection{Conexión de red}
Con el fin de minimizar la incertidumbre temporal (\emph{jitter}) y la latencia en el canal de comunicación, se recomienda establecer una conexión \textbf{directa} (Ethernet--RJ45) entre el \emph{Brain} OPTO22 y el computador que ejecuta la simulación. En lo posible, se debe evitar el uso de routers o switches intermedios, ya que introducen colas, procesamiento adicional y potenciales retransmisiones, aumentando la variabilidad del retardo entre paquetes.

Adicionalmente, se sugiere que el computador utilizado para la simulación opere en una red dedicada y, de ser viable, \textbf{sin acceso a Internet} durante la ejecución. Esto reduce tráfico de fondo (actualizaciones, sincronizaciones, servicios en segundo plano) y contribuye a mantener un canal de comunicación más estable, consistente y reproducible para el lazo de control en tiempo real.
